% संपूर्ण मराठी ग्रंथातील प्रकरण 49: मोडल क्रमवर्ती कलन
% हा संपादनयोग्य स्रोतखंड आहे. संकलनासाठी संपूर्ण स्रोतसंग्रहातील
% अक्षररूपे, पूर्वभाग, चिन्हव्याख्या आणि आकृती-स्रोत आवश्यक आहेत.
% मूळ: Open Logic Project; मजकूर आणि रूपांतर CC BY 4.0.
% यंत्रानुवाद, दुरुस्ती आणि तपासणी: OpenAI Codex —
% GPT-5.6 Sol आणि GPT-6 Sol, दोन्ही Ultra effort.
% दुरुस्ती-पुनरावलोकन: GPT-6.1 Sol, Ultra effort.
\chapter{मोडल क्रमवर्ती कलन}\label{nml:seq:chap}










\section{प्रस्तावना}\label{nml:seq:int:sec}

विधानीय तर्कशास्त्राच्या क्रमवर्ती कलनात
~$\Box$ आणि~
$\Diamond$ यांच्याशी संबंधित
अतिरिक्त नियम घालून त्याचा विस्तार करता येतो.
उदाहरणार्थ,~$\Log{K}$ साठी \Log{LK} बरोबर पुढील
नियम मिळतात:
\[    \Axiom$\Gamma \fCenter \Delta, \varphi$
    \RightLabel{$\Box$}
    \UnaryInf$\Box\Gamma \fCenter \Diamond\Delta, \Box\varphi$
    \DisplayProof
    \qquad
    \Axiom$\varphi, \Gamma \fCenter \Delta$
    \RightLabel{$\Diamond$}
    \UnaryInf$\Diamond\varphi, \Box\Gamma \fCenter \Diamond\Delta$
    \DisplayProof
\]
~$\Log{K}$ च्या विस्तारांसाठीही आणखी नियम
जोडावे लागतात.

प्रत्येक मोडल तर्कशास्त्रासाठी अशा प्रकारचे
क्रमवर्ती कलन उपलब्ध असेलच असे नाही. अर्थदृष्ट्या
साध्या असलेल्या~$\Log{S5}$ साठीसुद्धा (त्याची
व्याख्या प्राप्यता-संबंध मुळीच न वापरता करता येते),
मूळ मसुद्यात~$\Log{LK}$ पासून मिळणारे आणि~\Cut
नियमाशिवाय संपूर्ण असणारे कलन ज्ञात नसल्याचे म्हटले आहे.
हे सध्याच्या सर्व क्रमवर्ती कलनांविषयीचे विधान समजू नये.
उदाहरणार्थ, \href{https://arxiv.org/abs/1711.04634v3}{Aghaei
आणि Mohammadi यांच्या 2018 मधील लेखात} $\Log{G3c}$ चा
विस्तार असलेले $\Log{G3S5}$ दिले आहे. कलम~3 मधील नियम
आणि प्रमेये~5.1 व~5.2 त्यातील कटची अनुमतता आणि संपूर्णता
देतात; सिद्धतेसाठी त्याच कलनाचे अर्धविराम वापरणारे रूप घेतले आहे.
या प्रकरणात पुढे दिलेल्या विशिष्ट $\Log{LK}$-विस्तारासाठी
कटची गरज असल्याचा निष्कर्ष वेगळा आहे.
$\Log{S5}$ साठी कट-मुक्त संपूर्ण \emph{हायपरक्रमवर्ती}
कलनही आहे.












\section{\Log{K} साठी नियम}\label{nml:seq:rul:sec}

नेहमीच्या विधानीय संयोजकांसाठीचे नियम हे
नेहमीच्या क्रमवर्ती कलनातील~$\Log{LK}$ नियमांसारखेच
आहेत. स्वयंसिद्धकेही तीच आहेत: $\varphi \Sequent \varphi$
या रूपातील कोणतीही क्रमवर्ती स्वयंसिद्धक असते.

मोडल संकारक~$\Box$
आणि~$\Diamond$ यांच्यासाठी
पुढील अतिरिक्त नियम
आहेत:
\[    \Axiom$\Gamma \fCenter \Delta, \varphi$
    \RightLabel{$\Box$}
    \UnaryInf$\Box\Gamma \fCenter \Diamond\Delta, \Box\varphi$
    \DisplayProof
    \qquad
    \Axiom$\varphi, \Gamma \fCenter \Delta$
    \RightLabel{$\Diamond$}
    \UnaryInf$\Diamond\varphi, \Box\Gamma \fCenter \Diamond\Delta$
    \DisplayProof
\]
येथे $\Box\Gamma$ म्हणजे~$\Gamma$
मधील प्रत्येक सूत्राच्या आधी~$\Box$
लावून~$\Gamma$ पासून मिळणारा सूत्रांचा अनुक्रम;
आणि $\Diamond\Delta$ म्हणजे~$\Delta$
मधील प्रत्येक सूत्राच्या आधी~$\Diamond$
लावून~$\Delta$ पासून मिळणारा सूत्रांचा अनुक्रम.


$\Gamma$~आणि~$\Delta$
रिकामे असू शकतात; अशा वेळी निष्कर्ष-क्रमवर्तीतील
त्यांचे अनुक्रमे
$\Box\Gamma$~आणि~$\Diamond\Delta$
हे भागही रिकामे असतात.

एका सूत्र~$\varphi$ पुरतेच
उजव्या बाजूला~$\Box$ जोडणे आणि
डाव्या बाजूला~$\Diamond$ जोडणे
मर्यादित ठेवणे आवश्यक आहे. जर
उजवीकडे कोणत्याही संख्येने सूत्रे ना~$\Box$
जोडण्याची किंवा डावीकडे
कोणत्याही संख्येने सूत्रे ना~$\Diamond$ जोडण्याची
मुभा असती, तर पुढील निष्पन्न करता आले असते:
  \[      \Axiom$\varphi \fCenter \varphi$
      \RightLabel{\RightR{\lnot}}
      \UnaryInf$\fCenter \varphi, \lnot \varphi$
      \RightLabel{$\Box*$}
      \UnaryInf$\fCenter \Box\varphi, \Box\lnot \varphi$
      \doubleLine
      \RightLabel{\RightR{\lor}}
      \UnaryInf$\fCenter \Box\varphi \lor \Box \lnot \varphi$
      \DisplayProof
    \qquad
      \Axiom$\varphi \fCenter \varphi$
      \RightLabel{\LeftR{\lnot}}
      \UnaryInf$\lnot \varphi, \varphi \fCenter$
      \RightLabel{$\Diamond*$}
      \UnaryInf$\Diamond\lnot \varphi,\Diamond\varphi \fCenter $
      \RightLabel{\RightR{\lnot}}
      \UnaryInf$\Diamond\varphi \fCenter \lnot\Diamond\lnot \varphi$
      \RightLabel{\RightR{\lif}}
      \UnaryInf$ \fCenter \Diamond\varphi \lif \lnot\Diamond\lnot \varphi$
      \DisplayProof
  \]
मात्र $\Box\varphi \lor \Box \lnot \varphi$ आणि
$\Diamond\varphi \lif \lnot\Diamond\lnot \varphi$ यांपैकी कोणतेही सूत्र~$\Log{K}$ मध्ये
वैध नाही.

जर आधारविधानात~$\varphi$ व्यतिरिक्त बाजूची सूत्रेही
असू दिली, आणि $\Box$ नियमाने उजवीकडे
फक्त~$\varphi$ ला~$\Box$ जोडू दिला किंवा
$\Diamond$ नियमाने डावीकडे
फक्त~$\varphi$ ला~$\Diamond$ जोडू दिला (पण बाजूच्या
सूत्रांवर कोणतीही क्रिया केली नाही), तर पुढील
निष्पन्न करता आले असते:
\[    \Axiom$\varphi \fCenter \varphi$
    \RightLabel{\RightR{\lnot}}
    \UnaryInf$\fCenter \varphi, \lnot \varphi$
    \RightLabel{\RightR{\Exchange}}
    \UnaryInf$\fCenter \lnot \varphi, \varphi$
    \RightLabel{$\Box*$}
    \UnaryInf$\fCenter \lnot\varphi, \Box \varphi$
    \doubleLine
    \RightLabel{\RightR{\lor}}
    \UnaryInf$\fCenter \lnot\varphi \lor \Box \varphi$
    \DisplayProof
  \qquad
    \Axiom$\varphi \fCenter \varphi$
    \RightLabel{\LeftR{\lnot}}
    \UnaryInf$\lnot \varphi, \varphi \fCenter$
    \RightLabel{$\Diamond*$}
    \UnaryInf$\Diamond\lnot \varphi, \varphi \fCenter $
    \RightLabel{\RightR{\lnot}}
    \UnaryInf$\varphi \fCenter \lnot\Diamond\lnot \varphi$
    \RightLabel{\RightR{\lif}}
    \UnaryInf$ \fCenter \varphi \lif \lnot\Diamond\lnot \varphi$
    \DisplayProof
\]
मात्र $\lnot\varphi \lor \Box \varphi$ (जे $\varphi
\lif \Box\varphi$ शी सममूल्य आहे) आणि $\varphi \lif \lnot\Diamond\lnot \varphi$
 यांपैकी कोणतेही
सूत्र~$\Log{K}$ मध्ये वैध नाही.













\section{\Log{K} साठी क्रमवर्ती निष्पत्ती}\label{nml:seq:prk:sec}



\begin{ex}
  $\Proves (\Box\varphi \land \Box\psi)
  \lif \Box (\varphi \land \psi)$ हे दाखवणारी क्रमवर्ती कलनातील
  निष्पत्ती पुढे दिली आहे.
  \begin{prooftree}
    \Axiom$\varphi \fCenter \varphi$
    \doubleLine
    \UnaryInf$\psi, \varphi \fCenter \varphi$
    \Axiom$\psi \fCenter \psi$
    \doubleLine
    \UnaryInf$\psi, \varphi \fCenter \psi$
    \RightLabel{\RightR{\land}}
    \BinaryInf$\psi, \varphi \fCenter \varphi \land \psi$
    \RightLabel{$\Box$}
    \UnaryInf$\Box\psi, \Box\varphi \fCenter \Box
    (\varphi \land \psi)$
    \RightLabel{\LeftR{\land}}
    \UnaryInf$\Box\varphi \land \Box\psi, \Box\varphi \fCenter \Box
    (\varphi \land \psi)$
    \RightLabel{\LeftR{\Exchange}}
    \UnaryInf$\Box\varphi, \Box\varphi \land \Box\psi \fCenter \Box
    (\varphi \land \psi)$
    \RightLabel{\LeftR{\land}}
    \UnaryInf$\Box\varphi \land \Box\psi, \Box\varphi \land \Box\psi \fCenter \Box (\varphi \land \psi)$
    \RightLabel{\LeftR{\Contraction}}
    \UnaryInf$\Box\varphi \land \Box\psi \fCenter \Box (\varphi \land \psi)$
    \RightLabel{\RightR{\lif}}
    \UnaryInf$\fCenter (\Box\varphi \land \Box\psi)
    \lif \Box (\varphi \land \psi)$
  \end{prooftree}
\end{ex}



  \begin{ex}
    $\Proves \Diamond(\varphi
    \lor \psi) \lif (\Diamond \varphi \lor \Diamond\psi)$ हे दाखवणारी
    क्रमवर्ती कलनातील निष्पत्ती पुढे दिली आहे.
    \begin{prooftree}
      \Axiom$\varphi \fCenter \varphi$
      \doubleLine
      \UnaryInf$\varphi \fCenter \varphi, \psi$
      \Axiom$\psi \fCenter \psi$
      \doubleLine
      \UnaryInf$\psi \fCenter \varphi, \psi$
      \RightLabel{\LeftR{\lor}}
      \BinaryInf$\varphi \lor \psi \fCenter \varphi, \psi$
      \RightLabel{$\Diamond$}
      \UnaryInf$\Diamond(\varphi \lor \psi) \fCenter
      \Diamond \varphi, \Diamond\psi$
      \RightLabel{\RightR{\lor}}
      \UnaryInf$\Diamond(\varphi \lor \psi) \fCenter
      \Diamond \varphi, \Diamond \varphi \lor \Diamond\psi$
      \RightLabel{\RightR{\Exchange}}
      \UnaryInf$\Diamond(\varphi \lor \psi) \fCenter \Diamond \varphi \lor \Diamond\psi, \Diamond \varphi$
      \RightLabel{\RightR{\lor}}
      \UnaryInf$\Diamond(\varphi \lor \psi) \fCenter
      \Diamond \varphi \lor \Diamond\psi, \Diamond \varphi \lor \Diamond\psi$
      \RightLabel{\RightR{\Contraction}}
      \UnaryInf$\Diamond(\varphi
      \lor \psi) \fCenter \Diamond \varphi \lor \Diamond\psi$
      \RightLabel{\RightR{\lif}}
      \UnaryInf$\fCenter \Diamond(\varphi
      \lor \psi) \lif (\Diamond \varphi \lor \Diamond\psi)$
    \end{prooftree}
  \end{ex}



  येथे~\Dual ची निष्पत्ती दिली आहे.
  या सिद्धतेत $\psi \liff \chi$ हे
  $(\psi \lif \chi) \land (\chi \lif \psi)$ चे संक्षिप्त रूप घेतले आहे;
  म्हणून शेवटची पायरी दोन सशर्त सूत्रांचा संयोग बनवते.
  \begin{prooftree}
    \Axiom$\varphi \fCenter \varphi$
    \RightLabel{\LeftR{\lnot}}
    \UnaryInf$\lnot \varphi, \varphi \fCenter $
    \RightLabel{$\Diamond$}
    \UnaryInf$\Diamond\lnot \varphi, \Box \varphi \fCenter $
    \RightLabel{\RightR{\lnot}}
    \UnaryInf$\Box \varphi \fCenter \lnot\Diamond\lnot \varphi$
    \RightLabel{\RightR{\lif}}
    \UnaryInf$\fCenter \Box \varphi \lif \lnot\Diamond\lnot \varphi$
    \Axiom$\varphi \fCenter \varphi$
    \RightLabel{\RightR{\lnot}}
    \UnaryInf$\fCenter \varphi, \lnot \varphi $
    \RightLabel{\RightR{\Exchange}}
    \UnaryInf$\fCenter \lnot \varphi, \varphi $
    \RightLabel{$\Box$}
    \UnaryInf$\fCenter \Diamond\lnot \varphi, \Box \varphi$
    \RightLabel{\RightR{\Exchange}}
    \UnaryInf$\fCenter \Box \varphi, \Diamond\lnot \varphi$
    \RightLabel{\RightR{\lnot}}
    \UnaryInf$\lnot\Diamond\lnot \varphi \fCenter \Box \varphi$
    \RightLabel{\RightR{\lif}}
    \UnaryInf$\fCenter \lnot \Diamond \lnot \varphi \lif \Box \varphi$
    \RightLabel{\RightR{\land}}
    \BinaryInf$\fCenter \Box \varphi \liff \lnot\Diamond\lnot \varphi$
  \end{prooftree}








\begin{prob}
  पुढील सूत्रे साठी~$\Log{K}$ मध्ये
  क्रमवर्ती कलनातील सिद्धता शोधा:
  \begin{enumerate}
    \item $\Box \lnot p \lif \Box(p \lif q)$
    \item $(\Box p \lor \Box q) \lif \Box(p \lor q)$
    \item $\Diamond p \lif \Diamond(p \lor q)$
    \item $\Box(p \land q) \lif \Box p$
  \end{enumerate}
\end{prob}













\section{इतर प्राप्यता-संबंधांसाठी नियम}\label{nml:seq:mru:sec}

विशिष्ट प्राप्यता-संबंधांनी निर्धारित होणाऱ्या
तर्कशास्त्रांचा विचार करण्यासाठी आपण
\ref{nml:seq:mru:tab:more-rules} मधील अतिरिक्त नियम पाहू.


\begin{table}
  \begin{center}
    \def\arraystretch{3}
    \begin{tabular}{cc}
    \hline
      \Axiom$\varphi, \Gamma \fCenter \Delta$
      \RightLabel{T$\Box$}
      \UnaryInf$\Box\varphi, \Gamma \fCenter \Delta$
      \DisplayProof
    &
      \Axiom$\Gamma \fCenter \Delta, \varphi$
      \RightLabel{T$\Diamond$}
      \UnaryInf$\Gamma \fCenter \Delta, \Diamond\varphi$
      \DisplayProof
    \\[1ex]
    \multicolumn{2}{c}{
      \Axiom$\Gamma \fCenter \Delta$
      \RightLabel{D}
      \UnaryInf$\Box\Gamma \fCenter \Diamond\Delta$
      \DisplayProof}
    \\[1ex]
      \Axiom$\Gamma, \Diamond\Pi \fCenter \Box\Delta, \Lambda, \varphi$
      \RightLabel{B$\Box$}
      \UnaryInf$\Box\Gamma, \Pi \fCenter \Delta, \Diamond\Lambda, \Box\varphi$
      \DisplayProof
      &
      \Axiom$\varphi, \Diamond\Gamma, \Pi \fCenter \Box\Lambda, \Delta$
      \RightLabel{B$\Diamond$}
      \UnaryInf$\Diamond\varphi, \Gamma, \Box\Pi \fCenter \Lambda, \Diamond\Delta$
      \DisplayProof
    \\[1ex]
      \Axiom$\Box\Gamma \fCenter \Diamond\Delta, \varphi$
      \RightLabel{4$\Box$}
      \UnaryInf$\Box\Gamma \fCenter \Diamond\Delta, \Box\varphi$
      \DisplayProof
      &
      \Axiom$\varphi, \Box\Gamma \fCenter \Diamond\Delta$
      \RightLabel{4$\Diamond$}
      \UnaryInf$\Diamond\varphi, \Box\Gamma \fCenter \Diamond\Delta$
      \DisplayProof
    \\[1ex]
      \Axiom$\Box\Gamma, \Diamond\Pi \fCenter \Box\Delta,
      \Diamond\Lambda, \varphi$
      \RightLabel{5$\Box$}
      \UnaryInf$\Box\Gamma, \Diamond\Pi \fCenter \Box\Delta,
      \Diamond\Lambda, \Box\varphi$
      \DisplayProof
&
      \Axiom$\varphi, \Diamond\Gamma, \Box\Pi \fCenter \Diamond\Delta, \Box\Lambda$
      \RightLabel{5$\Diamond$}
      \UnaryInf$\Diamond\varphi,\Diamond\Gamma,\Box\Pi \fCenter \Diamond\Delta,\Box\Lambda$
      \DisplayProof
    \\[1ex]
    \hline
    \end{tabular}
  \end{center}
  \caption{अधिक मोडल नियम.}
  \label{nml:seq:mru:tab:more-rules}
\end{table}
हे नियम जोडल्यावर \ref{nml:seq:mru:tab:logics-rules} मध्ये
दिलेल्या तर्कशास्त्रांसाठी निर्दोष आणि संपूर्ण
पद्धती मिळतात.
\begin{table}
  \begin{center}
    \begin{tabular}{lll}
      \hline
      तर्कशास्त्र & $R$ हा \dots आहे & नियम\\
      \hline
      $\Log{T} = \Log{KT}$ & परावर्ती & $\Box$,
      $\Diamond$,
      T$\Box$
      ,
      T$\Diamond$
      \\ \hline
      $\Log{D} = \Log{KD}$ & सर्वत्र उत्तरवर्ती & $\Box$,
      $\Diamond$,
        D
      \\ \hline
      $\Log{K4}$ & संक्रमक & $\Box$,
      $\Diamond$,
      4$\Box$
      ,
      4$\Diamond$
      \\ \hline
      $\Log{B} = \Log{KTB}$ & परावर्ती, & $\Box$,
      $\Diamond$,
      T$\Box$
      ,
      T$\Diamond$\\
      & सममित &
      B$\Box$
      ,
      B$\Diamond$
      \\ \hline
      $\Log{S4} = \Log{KT4}$ & परावर्ती, & $\Box$,
      $\Diamond$,
      T$\Box$
      ,
      T$\Diamond$\\
      & संक्रमक &
      4$\Box$
      ,
      4$\Diamond$
      \\ \hline
      $\Log{S5} = \Log{KT5}$ & परावर्ती, & $\Box$,
      $\Diamond$,
      T$\Box$
      ,
      T$\Diamond$\\
      & संक्रमक, &
      5$\Box$
      ,
      5$\Diamond$\\
      &  युक्लिडी &
      \\ \hline
    \end{tabular}
  \end{center}
  \caption{विविध मोडल तर्कशास्त्रांसाठी क्रमवर्ती नियम.}
  \label{nml:seq:mru:tab:logics-rules}
\end{table}
\begin{ex}
  $\Log{K4} \Proves \Ax{4}$, म्हणजे
  $\Box\varphi \lif \Box\Box\varphi$, हे दाखवणारी क्रमवर्ती
  निष्पत्ती पुढे दिली आहे.
  \begin{prooftree}
    \Axiom$\Box\varphi \fCenter \Box\varphi$
    \RightLabel{4$\Box$}
    \UnaryInf$\Box\varphi \fCenter \Box\Box\varphi$
    \RightLabel{\RightR{\lif}}
    \UnaryInf$\fCenter \Box\varphi \lif \Box\Box\varphi$
  \end{prooftree}
\end{ex}
\begin{ex}
  $\Log{S5} \Proves \Ax{5}$, म्हणजे
  $\Diamond\varphi \lif \Box\Diamond\varphi$, हे दाखवणारी
  क्रमवर्ती निष्पत्ती पुढे दिली आहे.
  \begin{prooftree}
    \Axiom$\Diamond\varphi \fCenter \Diamond\varphi$
    \RightLabel{5$\Box$}
    \UnaryInf$\Diamond\varphi \fCenter \Box\Diamond\varphi$
    \RightLabel{\RightR{\lif}}
    \UnaryInf$\fCenter \Diamond\varphi \lif \Box\Diamond\varphi$
  \end{prooftree}
\end{ex}
\begin{ex}
या विभागात दिलेले \Log{S5} साठीचे क्रमवर्ती कलन \Cut{} नियमाशिवाय
संपूर्ण नाही; उदा., $\Diamond\Box \varphi \lif \varphi$ हे~$\Log{S5}$ मध्ये वैध असलेले
सूत्र~\Cut शिवाय सिद्ध करता येत नाही. पुढे~\Cut
वापरणारी निष्पत्ती दिली आहे:
\begin{prooftree}
\Axiom$\Box \varphi \fCenter \Box\varphi$
\RightLabel{5$\Diamond$}
\UnaryInf$\Diamond\Box\varphi \fCenter \Box\varphi$
\Axiom$\varphi \fCenter \varphi$
\RightLabel{T$\Box$}
\UnaryInf$\Box\varphi \fCenter \varphi$
\RightLabel{\Cut}
\BinaryInf$\Diamond\Box\varphi \fCenter \varphi$
\RightLabel{\RightR{\lif}}
\UnaryInf$\fCenter \Diamond\Box\varphi \lif \varphi$
\end{prooftree}
\end{ex}
\begin{prob}
पुढील निष्कर्ष दाखवणाऱ्या क्रमवर्ती निष्पत्त्या द्या:
  \begin{enumerate}
  \item $\Log{KT5} \Proves \Ax{B}$;
  \item $\Log{KT5} \Proves \Ax{4}$;
  \item $\Log{KDB4} \Proves \Ax{T}$;
  \item $\Log{KB4} \Proves \Ax{5}$;
  \item $\Log{KB5} \Proves \Ax{4}$;
  \item $\Log{KT} \Proves \Ax{D}$.
  \end{enumerate}
\end{prob}
\part{उपयोजित मोडल तर्कशास्त्र}\label{aml:part}
\begin{editorial}
या भागात कालिक आणि ज्ञानविषयक तर्कशास्त्रांसारख्या
मोडल तर्कशास्त्राच्या काही उपयोजनांवरील प्रायोगिक
मसुदा साहित्य आहे.
\end{editorial}
